\section{Transfer from the sphere to the cube}
\label{sec:orientation}

We now transfer the spherical conclusions to the Ising cube.
Conditional on the spectrum, GOE eigenvectors are Haar oriented.
Averaging over that orientation turns a single cube point into a
uniform sphere point. To use this first-moment identity for normalized
Gibbs probabilities, we must also control the random Ising partition
sum. A relative second-moment estimate supplies that control and then
transfers both small projections and the squared overlap.

The first-moment Haar-averaging identity is the fixed-spectrum spherical
comparison of Comets \cite[Equations~(2.1)--(2.2)]{comets1996}.
The fixed-spectrum Haar second-moment setup and its two-replica
reduction have antecedents in Bhattacharya and Sen
\cite[Section~3.2.1]{bhattacharyaSen2016}. Du and Huang
\cite[Theorem~1.6]{duHuang2026FreeEnergy} prove an unconditional shared-GOE
mean-square partition-ratio error of order $O(n^{-1/3})$ at criticality.
Here we prove the
conditional comparison uniformly on the particular spectral events of
Section~\ref{sec:spectral}. The central lattice comparison uses outward
bins and monotonicity; it needs no derivative bound for the reweighted
pair integral.

Throughout this section, let $L_n\to\infty$ with $L_n\le\log n$ and
$s_0=L_n n^{-2/3}$.  Write $\mathcal G_n$ for the eigenvalue-only good event in
Proposition~\ref{prop:spectral}.  We initially fix a deterministic sequence of
spectra in these events and use the notation
$\Lambda=\diag(\lambda_1,\ldots,\lambda_n)$, $S(a)$, $B(a)$, $z_a$, and
$\sigma_\Lambda$ from Section~\ref{sec:spherical}.  The error bounds supplied by
Proposition~\ref{prop:spectral} are deterministic on $\mathcal G_n$ and hold there
for every fixed finite upper resolvent parameter.

\subsection{Haar averaging and the two overlap laws}

Let $O$ be Haar distributed on the orthogonal group, independent of the
ordered spectrum, and put $W=O\Lambda O^T$.  Orthogonal invariance of GOE
justifies this representation: Haar averaging a test function of a symmetric
matrix depends only on its eigenvalues, and gives the invariant probability
measure on their conjugacy orbit.  We write $\E_O$ and $\Prob_O$ for
expectation and probability over $O$ with $\Lambda$ fixed.  Define
\begin{equation}
 I=2^{-n}\sum_{x\in\{-1,1\}^n}
       \exp\!\left(\frac{x^T O\Lambda O^T x}{2}\right).
 \label{eq:orientation-ising-partition}
\end{equation}

For every fixed cube point $x$, the vector $O^Tx$ is uniform on the
sphere of radius $\sqrt n$. Averaging each summand therefore gives the
exact first-moment identity
\[
 \E_O I=S(1).
\]
To control the normalization of the Gibbs law, we will compare the
second moment of $I$ with $S(1)^2$. Both quantities can be expressed
using the same pair integral; only their overlap laws differ.

Let
\[
 \mathcal Q_n=\{-1+2j/n:0\le j\le n\},
 \qquad
 p_I(q)=2^{-n}\binom{n}{n(1+q)/2}\quad(q\in\mathcal Q_n).
\]
This is the overlap law of two independent uniform cube points.

Let $(U,V)$ be a uniformly oriented orthogonal pair, each vector having norm
$\sqrt n$, and define, for $-1\le q\le1$,
\begin{equation}
 M(q)=\E\exp\!\left(
       \frac{(1+q)U^T\Lambda U+(1-q)V^T\Lambda V}{2}\right).
 \label{eq:orientation-pair-moment}
\end{equation}
If cube points $x,y$ have overlap $q\in(-1,1)$, their normalized sum and
difference
\[
 \frac{x+y}{\sqrt{2(1+q)}},
 \qquad
 \frac{x-y}{\sqrt{2(1-q)}}
\]
are orthogonal and both have norm $\sqrt n$.  Rotating them by $O^T$ gives
the pair in \eqref{eq:orientation-pair-moment}.  At $q=1$ or $q=-1$, the
points satisfy $y=x$ or $y=-x$, respectively, and the combined energy is
$x^T W x$; accordingly $M(1)=M(-1)=S(2)$.  Thus the second-moment identity
holds at every lattice point:
\begin{equation}
 \E_O I^2=\sum_{q\in\mathcal Q_n}p_I(q)M(q).
 \label{eq:orientation-second-moment}
\end{equation}

The exchangeability of $U,V$ shows that $M$ is even and nondecreasing in
$|q|$.  More explicitly, if $A=U^T\Lambda U$ and $D=V^T\Lambda V$, then
\[
 M(q)=\E\!\left[
 e^{(A+D)/2}\cosh\!\left(\frac{q(A-D)}2\right)\right],
\]
whose integrand is nondecreasing in $|q|$.

For two independent uniform sphere points the overlap has density
\[
 f_S(q)=c_n(1-q^2)^{(n-3)/2},
 \qquad
 c_n=\frac{\Gamma(n/2)}{\sqrt\pi\,\Gamma((n-1)/2)}
 \quad(-1<q<1).
\]
To see this, fix the first point by rotation and integrate the remaining
$n-1$ coordinates of the second point; the resulting one-coordinate
density is proportional to $(1-q^2)^{(n-3)/2}$, with the displayed
normalization given by the beta integral.  Conditional on this overlap,
the normalized sum and difference are again the orthogonal pair above.
Independence of the two spherical points consequently gives the exact identity
\begin{equation}
 \int_{-1}^1 M(q)f_S(q)\,dq=S(1)^2.
 \label{eq:orientation-spherical-second-moment}
\end{equation}

The moment identities reduce concentration to a comparison of the two
overlap laws. Because $\E_O I=S(1)$, it is enough to prove
\[
 \sum_{q\in\mathcal Q_n}p_I(q)M(q)
 \le (1+o(1))\int_{-1}^1 M(q)f_S(q)\,dq
\]
with a deterministic $o(1)$ uniform over $\mathcal G_n$. We will use a
variational estimate for large overlaps and a direct comparison of the
unweighted overlap laws near zero.

First consider the size of the pair integral. By evenness it suffices
to take $q\ge0$. Conditional on $U$ in
\eqref{eq:orientation-pair-moment}, the vector $V$ is uniform on the
radius-$\sqrt n$ sphere in $U^\perp$. The codimension-one bound of
Lemma~\ref{lem:partition-bounds} controls its integral at parameter
$1-q$, uniformly in $U$. The full-space bound then controls the
remaining integral at parameter $1+q$. Dividing by the lower bound for
$S(1)^2$ from the same lemma gives
\begin{equation}
 \frac{M(q)}{S(1)^2}
 \le Cn^2\exp\!\left(n\bigl[B(1+|q|)+B(1-|q|)-2B(1)\bigr]\right).
 \label{eq:orientation-pair-bound}
\end{equation}
This remains valid at $|q|=1$, since the integral with parameter zero
equals one.

The cube overlap mass satisfies $p_I(q)\le e^{-nq^2/2}$. Indeed, for a
mean of $n$ independent uniform signs, the exponential moment bound
$\cosh t\le e^{t^2/2}$ follows by integrating
$(\log\cosh t)''\le1$ twice from zero. Optimizing the resulting Chernoff
bound and using symmetry gives the stated bound on each lattice mass.
Thus a margin below $q^2/2$ in the exponent of
\eqref{eq:orientation-pair-bound} will suppress large overlaps.

We split the overlap range at
\[
 q_0=n^{-7/24}.
\]
The proof uses three scale relations:
\[
 nq_0^3=n^{1/8}\longrightarrow\infty,\qquad
 nq_0^4=n^{-1/6}\longrightarrow0,\qquad
 \frac{s_0}{q_0^2}=L_n n^{-1/12}\longrightarrow0.
\]
The first turns a cubic variational margin into a negligible tail. The
second permits a relative local comparison of the unweighted cube and
sphere overlap laws. The third absorbs the spectral error in the cubic
estimate; it uses $L_n\le\log n$.

\subsection{The variational deficit}

The next lemma supplies the margin needed in
\eqref{eq:orientation-pair-bound} for $|q|\ge q_0$.

\begin{lemma}[Cubic deficit]
\label{lem:cubic-deficit}
There is a universal constant $c>0$ such that, for all sufficiently large $n$,
uniformly for spectra in $\mathcal G_n$ and $q_0\le q\le1$,
\begin{equation}
 B(1+q)+B(1-q)-2B(1)
 \le \frac{q^2}{2}-cq^3.
 \label{eq:cubic-deficit}
\end{equation}
One may take $c=1/24$.
\end{lemma}

\begin{proof}
We estimate the derivative of the spherical variational value on
both sides of the critical parameter $a=1$. Close to that parameter
the resolvent estimates give a quantitative margin; away from it a
bounded-convergence argument suffices. The saddle equation in
Lemma~\ref{lem:spherical-identity} and implicit differentiation show
that $z_a$ is continuously differentiable in $a$.
Differentiating the objective at its minimizer gives
\begin{equation}
 B'(a)=\frac1{2n}\sum_{i=1}^n
       \frac{\lambda_i}{z_a-a\lambda_i},
 \qquad
 B'(a)=\frac{z_a-1}{2a}\quad(a>0).
 \label{eq:orientation-saddle-derivative}
\end{equation}
The first identity also holds at $a=0$, with a one-sided derivative at the
endpoint.  Since the positive weights $(z_a-a\lambda_i)^{-1}$ sum to $n$,
\[
 |B'(a)|\le\frac12\max_i|\lambda_i|\le\frac32
 \qquad(0\le a\le2).
\]
For $a>0$, put $r_a=z_a/a$ and
\[
 m_n(s)=\frac1n\sum_i(2+s-\lambda_i)^{-1}.
\]
The saddle equation is equivalently
\[
 \frac1n\sum_i(r_a-\lambda_i)^{-1}=a.
\]
At $s=s_0$, the spectral estimates give
$m_n(s_0)=1-\sqrt{s_0}+o(\sqrt{s_0})<1$.  Thus, for $a=1+u$ with
$0\le u\le1$, monotonicity of this resolvent gives $r_a<2+s_0$ and hence
\begin{equation}
 B'(1+u)
 \le \frac{1+u}{2}-\frac{u^2}{2(1+u)}+\frac{s_0}{2}.
 \label{eq:orientation-supercritical-derivative}
\end{equation}

The upper derivative estimate must be paired with a lower estimate
below $a=1$. The small-overlap range requires the following quantitative
comparison. If
$4\sqrt{s_0}\le u\le1/4$, set $t=u^2/(1-u)$.  Then $t\le1/12$ and
$0.9t\ge14.4s_0$.  The function $R$ from
Proposition~\ref{prop:spectral} satisfies
\[
 R(t)=1-u,\qquad
 -R'(s)=\frac{2+s}{2\sqrt{s(4+s)}}-\frac12
        \ge \frac{c_1}{\sqrt{s}}
        \quad(0<s\le1/12)
\]
for a universal $c_1>0$.  Integrating the last bound from $0.9t$ to $t$
therefore gives
\[
 R(0.9t)-(1-u)\ge c_2\sqrt{t}\ge c_2u.
\]
The error in $m_n(0.9t)=R(0.9t)+o(\sqrt{t})$ is uniform over the indicated
range.  It is smaller than this fixed positive margin for all sufficiently
large $n$.  Consequently $m_n(0.9t)>1-u$, so monotonicity implies
$r_{1-u}\ge2+0.9t$.  Substitution in
\eqref{eq:orientation-saddle-derivative} yields
\begin{equation}
 B'(1-u)\ge\frac{1-u}{2}-\frac{u^2}{20(1-u)}
 \qquad(4\sqrt{s_0}\le u\le1/4).
 \label{eq:orientation-subcritical-derivative}
\end{equation}
For $0\le u\le4\sqrt{s_0}$, use instead
$r_{1-u}>\lambda_1=2+o(s_0)$ to obtain
\[
 B'(1-u)
 \ge\frac{1-u}{2}-\frac{u^2}{2(1-u)}-o(s_0)
 \ge\frac{1-u}{2}-Cs_0.
\]

Write $F(q)=B(1+q)+B(1-q)-2B(1)$.  On $0\le u\le1/4$,
\[
 \frac1{2(1+u)}-\frac1{20(1-u)}\ge\frac13.
\]
Combining the preceding bounds, and absorbing $u^2\le16s_0$ on the short
interval $u\le4\sqrt{s_0}$ into the error, gives
\[
 F'(u)\le u-\frac{u^2}{3}+Cs_0
 \qquad(0\le u\le1/4).
\]
Since $F(0)=0$,
\begin{equation}
 F(q)\le\frac{q^2}{2}-\frac{q^3}{9}+Cs_0q
 \qquad(0\le q\le1/4).
 \label{eq:orientation-small-overlap-deficit}
\end{equation}
For $q\ge q_0$, the last error divided by $q^3$ is at most
\[
 C\frac{s_0}{q_0^2}=CL_n n^{-1/12}\longrightarrow0.
\]
Thus \eqref{eq:cubic-deficit} holds on $[q_0,1/4]$, even with $c=1/18$.

It remains to cover overlaps bounded away from zero. Here a
uniform additive error tending to zero is enough. Fix
$a\in(0,1)$. The identity
$R(a+a^{-1}-2)=a$, the strict monotonicity of $R$, and the spectral resolvent
estimate imply
\[
 r_a\longrightarrow a+a^{-1},
 \qquad B'(a)\longrightarrow a/2.
\]
Indeed, evaluating the empirical resolvent at
$a+a^{-1}\pm\delta$, with $\delta>0$ small enough that the lower point is
above $2$, traps $r_a$ between these two points for all large $n$.
The upper resolvent parameter needed here is finite for each fixed $a$.
No estimate uniform as $a\downarrow0$ is used.

The bound $|B'|\le3/2$ now gives, by dominated convergence,
\[
 \rho_n:=\int_0^1|B'(a)-a/2|\,da\longrightarrow0.
\]
Integrating \eqref{eq:orientation-supercritical-derivative} and bounding the
subcritical error by $\rho_n$ gives, uniformly for $0\le q\le1$,
\begin{align}
 F(q)
 &\le \frac{q^2}{2}
       -\int_0^q\frac{u^2}{2(1+u)}\,du+\rho_n+\frac{s_0}{2} \notag\\
 &\le \frac{q^2}{2}-\frac{q^3}{12}+\rho_n+\frac{s_0}{2}.
 \label{eq:orientation-macroscopic-deficit}
\end{align}
When $q\ge1/4$, the cubic term is bounded away from zero, so the last two
errors are absorbed with $c=1/24$.  This also covers $q=1$.

The conclusion must hold uniformly over the entire deterministic
spectral event. The estimates in the shrinking range already have
that uniformity.
The same is true of $\rho_n\to0$: otherwise one could select spectra from
$\mathcal G_n$ along a violating subsequence.  Their deterministic spectral
bounds would still imply the pointwise convergence and bounded domination
just proved, a contradiction.  Thus the two ranges give the claimed
uniform statement.
\end{proof}

\subsection{Concentration over the orientation}

\begin{proposition}[Conditional partition concentration]
\label{prop:orientation-concentration}
For spectra in $\mathcal G_n$,
\[
 \E_O I=S(1),
 \qquad
 \sup_{\Lambda\in\mathcal G_n}
 \E_O\!\left[\left(\frac{I}{S(1)}-1\right)^2\right]\longrightarrow0.
\]
In particular, under the GOE law, $I/S(1)\to1$ in probability.
\end{proposition}

\begin{proof}
The first-moment identity was established above. We now bound the exact
second moment by treating the two overlap regions separately.

\smallskip\noindent\emph{Large overlaps.}
Combining \eqref{eq:orientation-pair-bound}, the Chernoff bound
$p_I(q)\le e^{-nq^2/2}$, and Lemma~\ref{lem:cubic-deficit} gives
\[
 p_I(q)\frac{M(q)}{S(1)^2}
 \le Cn^2 e^{-cn|q|^3}
 \qquad(q\in\mathcal Q_n,\ |q|\ge q_0).
\]
There are at most $n+1$ lattice points. Since $nq_0^3=n^{1/8}$,
\begin{equation}
 \sum_{\substack{q\in\mathcal Q_n\\ |q|\ge q_0}}
       p_I(q)\frac{M(q)}{S(1)^2}
 \le Cn^3e^{-c n^{1/8}}=: \tau_n.
 \label{eq:orientation-tail}
\end{equation}
In particular, $\tau_n$ is smaller than every fixed negative power of $n$.

\smallskip\noindent\emph{Central overlaps.}
The remaining obstacle is that $M(q)$ may vary too rapidly for a
derivative-based Riemann-sum estimate. We compare the unweighted overlap
laws locally, then move each lattice mass outward, where $M$ is larger.
Let $\Delta=2/n$. Stirling's formula,
uniformly for $|q|<q_0$, gives
\begin{gather*}
 p_I(q)=
 \sqrt{\frac{2}{\pi n(1-q^2)}}\,
 e^{-nJ(q)}\bigl(1+O(n^{-1})\bigr),\\
 J(q)=\frac{(1+q)\log(1+q)+(1-q)\log(1-q)}2.
\end{gather*}
Here $J(q)=q^2/2+O(q^4)$.  Likewise
$c_n=\sqrt{n/(2\pi)}(1+O(n^{-1}))$ and
$\log(1-q^2)=-q^2+O(q^4)$.  Dividing the two expressions shows that
\begin{equation}
 \frac{p_I(q)}{\Delta f_S(q)}
 =\exp\!\bigl(O(n^{-1}+q^2+nq^4)\bigr)=1+o(1),
 \qquad |q|<q_0,
 \label{eq:orientation-local-density}
\end{equation}
with a uniform relative error, since $nq_0^4\to0$.

Assign to each retained lattice point the outward bin
\begin{equation}
 \mathcal B_q=
 \begin{cases}
 [q,q+\Delta],&q>0,\\
 [q-\Delta,q],&q<0,\\
 [-\Delta/2,\Delta/2],&q=0.
 \end{cases}
 \label{eq:orientation-bins}
\end{equation}
For even $n$, the lattice near zero is
$\ldots,-2\Delta,-\Delta,0,\Delta,2\Delta,\ldots$.
The centered zero bin leaves gaps between $\Delta/2$ and $\Delta$, and
between $-\Delta$ and $-\Delta/2$.
For odd $n$, the nearest points are $-\Delta/2,\Delta/2$, whose outward
bins leave the gap $(-\Delta/2,\Delta/2)$.
In both cases all bin interiors are disjoint.  Every $r\in\mathcal B_q$
satisfies $|r|\ge|q|$, and all these bins lie in
$[-q_0-\Delta,q_0+\Delta]\subset(-1,1)$ for large $n$.

On each such bin,
\[
 \left|\log\frac{f_S(r)}{f_S(q)}\right|
 \le Cn(q_0+\Delta)\Delta\le C(q_0+\Delta),
\]
because $(\log f_S)'(r)=-(n-3)r/(1-r^2)$.
It follows from \eqref{eq:orientation-local-density} that there is a
deterministic $\delta_n\to0$ such that, for every nonnegative even function
$h$ nondecreasing in $|q|$,
\begin{equation}
 p_I(q)h(q)
 \le (1+\delta_n)\int_{\mathcal B_q}h(r)f_S(r)\,dr
 \qquad(q\in\mathcal Q_n,\ |q|<q_0).
 \label{eq:orientation-bin-comparison}
\end{equation}
Indeed, the integral is at least
$\Delta h(q)f_S(q)e^{-C(q_0+\Delta)}$, whereas the lattice mass is at most
$(1+e_n)\Delta f_S(q)$ for a deterministic $e_n\to0$.
One can take $1+\delta_n=(1+e_n)e^{C(q_0+\Delta)}$.
This estimate does not depend on how rapidly $h$ varies.
The outermost bin may extend beyond $q_0$, which causes no difficulty
because the comparison integral is over the full interval $[-1,1]$.

\smallskip\noindent\emph{Combining the two regions.}
Apply \eqref{eq:orientation-bin-comparison} with $h=M$. Because the bin
interiors are disjoint, their sum is bounded by the full spherical integral.  Together with \eqref{eq:orientation-tail} and
\eqref{eq:orientation-spherical-second-moment}, it yields
\[
 \frac{\E_O I^2}{S(1)^2}\le1+\delta_n+\tau_n.
\]
The first-moment identity therefore gives
\begin{equation}
 \E_O\!\left[\left(\frac I{S(1)}-1\right)^2\right]
 \le\delta_n+\tau_n\longrightarrow0.
 \label{eq:orientation-relative-variance}
\end{equation}
The bin error is independent of the spectrum, and the tail estimate is
uniform on $\mathcal G_n$ by Lemma~\ref{lem:cubic-deficit}.  Thus the
relative variance bound is uniform there. The unconditional assertion
follows by Chebyshev's inequality and $\Prob(\mathcal G_n^c)\to0$.
\end{proof}

\subsection{Small projections and the squared overlap}

\begin{theorem}[Small projections in the Ising law]
\label{thm:static-small-ball}
Let $W$ be the augmented GOE matrix and let $v$ be a unit top eigenvector.
For every positive deterministic sequence $b_n\to0$,
\begin{equation}
 \mu\{x:|v^Tx|\le b_n n^{1/3}\}
 \longrightarrow0
 \quad\text{in probability}.
 \label{eq:static-small-ball}
\end{equation}
The probability here includes the auxiliary diagonal used to define $v$;
the Gibbs measure itself does not depend on that diagonal.
\end{theorem}

\begin{proof}
The restricted partition sum has a small first moment, while the
unrestricted sum is close to its spherical mean. We choose the auxiliary
spectral event after the prescribed threshold $b_n$: take $L_n\to\infty$
with $L_n\le\log n$ and $L_n^2b_n\to0$.
For example, for all sufficiently large $n$ one may take
$L_n=\min\{\log n,b_n^{-1/4}\}$.
In the representation $W=O\Lambda O^T$, take $v=Oe_1$, and define the
restricted normalized sum
\[
 I_b=2^{-n}\sum_x
      \mathbf 1_{\{|v^Tx|\le b_n n^{1/3}\}}e^{x^TWx/2}.
\]
For every fixed $x$, $O^Tx$ is uniform on the sphere and its first
coordinate is $v^Tx$.  Consequently Proposition~\ref{prop:spherical-small-ball}
gives, conditional on any spectrum in $\mathcal G_n$,
\begin{equation}
 \E_O\frac{I_b}{S(1)}
 =\sigma_\Lambda\{|U_1|\le b_n n^{1/3}\}
 \le CL_n^2b_n.
 \label{eq:orientation-restricted-first-moment}
\end{equation}
For every fixed $\varepsilon>0$, Markov's inequality and
Proposition~\ref{prop:orientation-concentration} imply
\[
 \Prob_O\!\left(\frac{I_b}{I}>\varepsilon\right)
 \le \Prob_O\!\left(\frac{I}{S(1)}<\frac12\right)
       +\frac{2}{\varepsilon}\E_O\frac{I_b}{S(1)}
 \longrightarrow0
\]
uniformly over the good spectra.  The ratio $I_b/I$ is the desired Gibbs
mass, so adding $\Prob(\mathcal G_n^c)$ proves the claim.
GOE has a simple spectrum almost surely: the set of matrices with repeated
eigenvalues is the zero set of the nonzero discriminant polynomial and
has Lebesgue measure zero, whereas GOE has a density.
Thus any choice of the unit top eigenvector differs from $Oe_1$ only by
sign, which leaves the asserted event unchanged.
\end{proof}

\begin{proposition}[Squared overlap]
\label{prop:overlap}
Let $Q_n=\E_{\mu^{\otimes2}}q(x,y)^2$.  Then $n^{2/3}Q_n$ is bounded in
probability.  Equivalently, for every positive deterministic sequence
$T_n\to\infty$,
\begin{equation}
 \Prob\{n^{2/3}Q_n>T_n\}\longrightarrow0.
 \label{eq:orientation-overlap-tail}
\end{equation}
This assertion concerns the original off-diagonal SK disorder.
\end{proposition}

\begin{proof}
To prove tightness rather than a bound with a fixed logarithmic
loss, keep $L_n$ free until the tail threshold has been specified.
The same second-moment comparison works with the nonnegative overlap
weight $q^2$. Expanding two Gibbs weights gives the exact identity
\begin{equation}
 \E_O\!\left[\left(\frac I{S(1)}\right)^2Q_n\right]
 =\frac1{S(1)^2}\sum_{q\in\mathcal Q_n}p_I(q)q^2M(q).
 \label{eq:orientation-weighted-second-moment}
\end{equation}
The function $q^2M(q)$ is nonnegative, even, and nondecreasing in $|q|$.
The central terms therefore satisfy the same bin comparison
\eqref{eq:orientation-bin-comparison}; the tails are at most $\tau_n$
because $q^2\le1$.  Moreover,
\[
 \frac1{S(1)^2}\int_{-1}^1q^2M(q)f_S(q)\,dq
 =\E_{\sigma_\Lambda^{\otimes2}}
       \left(\frac{U^TV}{n}\right)^2.
\]
Here $U,V$ are independent under the tilted spherical law.
Lemma~\ref{lem:spherical-covariance} bounds this expectation by
$Cn^{-2/3}L_n^4$.  Thus \eqref{eq:orientation-weighted-second-moment} gives
\begin{equation}
 \E_O\!\left[\left(\frac I{S(1)}\right)^2 n^{2/3}Q_n\right]
 \le CL_n^4+n^{2/3}\tau_n
 \le C'L_n^4
 \label{eq:orientation-weighted-bound}
\end{equation}
for all large $n$, uniformly on $\mathcal G_n$.

The weighted estimate is now available for every admissible $L_n$.
Given an arbitrary deterministic $T_n\to\infty$, choose
$L_n=\min\{\log n,T_n^{1/8}\}$ for all sufficiently large $n$.
Then $L_n\to\infty$ and $L_n^4/T_n\le T_n^{-1/2}\to0$.
Writing $\widehat I=I/S(1)$, conditional Markov's inequality gives
\[
 \Prob_O(n^{2/3}Q_n>T_n)
 \le \Prob_O(\widehat I<1/2)
       +\frac4{T_n}\E_O[\widehat I^{\,2} n^{2/3}Q_n]
 \le o(1)+\frac{CL_n^4}{T_n}.
\]
The first error is uniform on $\mathcal G_n$ by
Proposition~\ref{prop:orientation-concentration}.
Adding $\Prob(\mathcal G_n^c)\to0$ proves
\eqref{eq:orientation-overlap-tail}.

The random variable $n^{2/3}Q_n$ has one fixed law for each $n$ throughout
this argument: it is a function of the off-diagonal disorder alone.
Changing $L_n$ changes only the auxiliary good event and estimates used
to prove the assertion, not this law.  To verify explicitly that
\eqref{eq:orientation-overlap-tail} implies boundedness in probability,
suppose tightness failed.  Then there would be an $\varepsilon>0$ and
increasing indices $n_j$ with
\[
 \Prob(n_j^{2/3}Q_{n_j}>j)>\varepsilon.
\]
Define a deterministic staircase by $T_n=j$ for $n_j\le n<n_{j+1}$, with
an arbitrary positive initial segment.  This sequence tends to infinity
and contradicts \eqref{eq:orientation-overlap-tail} at $n=n_j$.
Conversely, tightness directly implies
\eqref{eq:orientation-overlap-tail} by comparing a divergent threshold
with any fixed threshold.  Finally, because the Gibbs law is invariant
under adding a diagonal to $W$, integrating out the auxiliary diagonal
preserves the asserted probabilities.
\end{proof}

\begin{corollary}[Covariance scale]
\label{cor:covariance-scale}
For every positive deterministic sequence $M_n\to\infty$,
\[
 \Prob\!\left\{
 \frac{n^{2/3}}{M_n}\le\|\Sigma\|\le M_n n^{2/3}
 \right\}\longrightarrow1.
\]
\end{corollary}

\begin{proof}
The upper and lower bounds use different outputs of the transfer:
squared overlap bounds the covariance norm from above, while the top
projection bounds it from below. Spin-flip symmetry gives $\mu[x]=0$,
so $\Sigma=\mu[xx^T]$ and
\begin{equation}
 \Tr\Sigma^2
 =\sum_{i,j}\mu[x_ix_j]^2
 =n^2\E_{\mu^{\otimes2}}q(x,y)^2
 =n^2Q_n.
 \label{eq:orientation-covariance-trace}
\end{equation}
Since $\Sigma$ is positive semidefinite,
$\|\Sigma\|^2\le\Tr\Sigma^2$.  The upper bound follows from
Proposition~\ref{prop:overlap} with threshold $T_n=M_n^2$.

For the lower bound use Theorem~\ref{thm:static-small-ball} with
$b_n=M_n^{-1/4}$ and put
$B_n=\mu\{|v^Tx|\le b_n n^{1/3}\}$.  Then $B_n\to0$ in probability and
\[
 \|\Sigma\|\ge v^T\Sigma v
 =\mu[(v^Tx)^2]
 \ge b_n^2 n^{2/3}(1-B_n).
\]
On $B_n\le1/2$, for all sufficiently large $n$ this is at least
$n^{2/3}/(2\sqrt{M_n})\ge n^{2/3}/M_n$.
Both bounds hold simultaneously with probability tending to one by a
union bound.  The covariance is unchanged by the auxiliary diagonal,
so these probabilities are also those of the original SK disorder.
\end{proof}
